\documentclass[tikz,border=6pt]{standalone}
\usepackage{ctex}
\usepackage{amsmath}
\usepackage{pgfplots}
\pgfplotsset{compat=1.18}
\usetikzlibrary{arrows.meta, decorations.pathmorphing}

\begin{document}
\begin{tikzpicture}

\begin{axis}[
  width=10cm, height=9cm,
  xmin=-0.2, xmax=4.5,
  ymin=-0.2, ymax=4.5,
  xlabel={$x_1$}, ylabel={$x_2$},
  xlabel style={font=\small},
  ylabel style={font=\small},
  axis lines=left,
  tick style={font=\tiny},
  xtick={0,1,2,3,4},
  ytick={0,1,2,3,4},
  clip=false,
]

% ── 可行域（阴影多边形）：x1>=0.3, x2>=0.3, x1+x2<=4.0 ────────────────────
\fill[blue!8]
  (axis cs:0.3, 0.3) --
  (axis cs:3.7, 0.3) --
  (axis cs:0.3, 3.7) -- cycle;

\draw[blue!50, thick]
  (axis cs:0.3, 0.3) -- (axis cs:3.7, 0.3)
  node[right, font=\tiny, blue!70] {$g_2(x)\!=\!0$};
\draw[blue!50, thick]
  (axis cs:0.3, 0.3) -- (axis cs:0.3, 3.7);
\node[font=\tiny, blue!70, anchor=east] at (axis cs:0.28, 2.0) {$g_1(x)\!=\!0$};
\draw[blue!50, thick]
  (axis cs:3.7, 0.3) -- (axis cs:0.3, 3.7);
\node[font=\tiny, blue!70, anchor=south west] at (axis cs:2.2, 1.9) {$g_3(x)\!=\!0$};

\node[blue!60, font=\small] at (axis cs:1.4, 1.6) {可行域 $\mathcal{F}$};

% ── 中心路径（从内部出发趋向最优解）─────────────────────────────────────────
% 最优解约在 (0.3, 0.3) 附近（左下角顶点）
% 参数路径：x*(t) = (0.3 + 1.2/t, 0.3 + 1.2/t) for t from 0.5 to large
% t→0: x*(t) → (2.7, 2.7) 附近（中心）；t→∞: 趋向最优解 (0.3, 0.3)

% ── 中心路径（Bezier 曲线）──────────────────────────────────────────────────
\draw[red!80!black, very thick, -{Stealth[length=7pt]}]
  (axis cs:2.8, 3.0)
  .. controls (axis cs:2.2, 2.5) and (axis cs:1.6, 2.0) ..
  (axis cs:1.2, 1.5)
  .. controls (axis cs:0.9, 1.2) and (axis cs:0.7, 0.9) ..
  (axis cs:0.55, 0.7)
  .. controls (axis cs:0.45, 0.55) and (axis cs:0.38, 0.42) ..
  (axis cs:0.33, 0.35);

% ── 中心路径上各点标注 ────────────────────────────────────────────────────────
\fill[red] (axis cs:2.8, 3.0) circle (2.5pt);
\node[font=\tiny, anchor=west, text=red!80!black] at (axis cs:2.85, 3.0)
  {$\mathbf{x}^*(t_1)$，$t_1$ 小};

\fill[red] (axis cs:1.2, 1.5) circle (2.5pt);
\node[font=\tiny, anchor=west, text=red!80!black] at (axis cs:1.25, 1.5)
  {$\mathbf{x}^*(t_2)$};

\fill[red] (axis cs:0.55, 0.7) circle (2.5pt);
\node[font=\tiny, anchor=west, text=red!80!black] at (axis cs:0.60, 0.7)
  {$\mathbf{x}^*(t_3)$};

% ── 最优解 ─────────────────────────────────────────────────────────────────────
\fill[black] (axis cs:0.3, 0.3) circle (3pt);
\node[font=\small, anchor=north east] at (axis cs:0.28, 0.28)
  {$\mathbf{x}^*$};

% ── 参数 t 增大方向箭头 ──────────────────────────────────────────────────────
\node[font=\small, red!80!black, align=center] at (axis cs:3.5, 3.8)
  {中心路径 $\mathcal{C}$};
\draw[-{Stealth}, red!50, thick]
  (axis cs:3.3, 3.7) -- (axis cs:3.0, 3.2);

% t 增大说明
\node[font=\footnotesize, align=left] at (axis cs:3.8, 2.0)
  {$t$ 增大\\$\rightarrow$\\趋近 $\mathbf{x}^*$};

% 标注方向弯头（从 2.8,3.0 到 0.3,0.3）
\draw[-{Stealth[length=5pt]}, gray, dashed]
  (axis cs:3.5, 2.1) to[bend right=20] (axis cs:2.8, 3.0);

% ── 等值线（目标函数 f 的轮廓线）──────────────────────────────────────────
% 示意：两条等值线（虚线椭圆）
\draw[gray!60, dashed, thick] (axis cs:0.3,0.3) ellipse [x radius=1.0cm, y radius=0.8cm];
\draw[gray!40, dashed] (axis cs:0.3,0.3) ellipse [x radius=1.8cm, y radius=1.5cm];
\node[font=\tiny, gray!70] at (axis cs:1.2, 0.7) {$f$ 等值线};

\end{axis}
\end{tikzpicture}
\end{document}
